\section{The Physical-AI-Fusion Benchmark}
\label{sec:benchmark}

\subsection{Scope and topic taxonomy}
The benchmark covers fusion physics with an ICF emphasis (Figure~\ref{fig:topics}): eight ICF subdomains---ignition \& gain, Rayleigh--Taylor/Richtmyer--Meshkov (RT/RM) instabilities, laser--plasma interactions (LPI), radiation hydrodynamics, equations of state \& opacity, implosion dynamics, target design, and diagnostics---plus magnetic confinement fusion (MCF), Z-pinch and pulsed power, plasma fundamentals, fusion engineering (materials, tritium breeding, blankets), and computational physics (numerical methods, classic validation problems). Topic selection follows the standard ICF curriculum (Atzeni~\cite{atzeni2004icf}; Lindl~\cite{lindl1998icf}; Drake~\cite{drake2018hedp}; Kruer~\cite{kruer1988lpi}; Zel'dovich \& Raizer~\cite{zeldovich1966shock}; Freidberg~\cite{freidberg2007fusion}) and the NRL Plasma Formulary~\cite{nrlformulary}.

\subsection{Task types and difficulty}
Problems are authored to be self-contained (all constants and unit conventions in the statement) and fall into five types: \textbf{numeric} (531; deterministic tolerance-banded grading, default relative tolerance 2\%), \textbf{concept} (135; single- and multi-select exact match, distractors engineered from common failure modes---unit-system slips, coefficient slips, regime misidentification), \textbf{derivation} (23; algebraic derivations with a standard symbolic answer in a declared symbol alphabet, graded by SymPy symbolic equivalence in the style of PHYSICS~\cite{feng2025physics}---simplification to zero, constant-factor proportionality, or sampled numeric equivalence), \textbf{code} (9; eight graded by containerized unit tests, one SWE-bench-style patch task graded by applying the response's unified diff and running FAIL\_TO\_PASS tests on a real FreeGS issue), and \textbf{figure} (1; table/distribution reading with tolerance). \textbf{All 699 tasks are graded deterministically}---no LLM judge is required anywhere in the pipeline. Difficulty spans undergrad (113), grad (373), and expert (213); expert tasks require multi-step chains (e.g., NIF energetics: laser$\to$X-ray conversion$\to$capsule absorption), cross-formula synthesis, or reading verified local simulation output. 27 problems derive from real GitHub issues of fusion simulation codes (DESC, OpenMC, WarpX, TORAX, BOUT++, FreeGS).

\subsection{Schema and sources}
Each task is one JSONL record (schema v1.0) carrying \texttt{task\_id}, \texttt{topic}, \texttt{type}, \texttt{difficulty}, \texttt{question}, a typed \texttt{answer}, a \texttt{grading} mode, provenance \texttt{metadata} (\texttt{source}, \texttt{source\_type}, \texttt{license}), and a \texttt{verification} block (methods, evidence, verifier, timestamp). Sources are public physics only: the NRL Plasma Formulary, textbook-standard formulas (quoted as formulas, never copied prose), LLNL/NIF public ignition data (2022-12-05: 2.05\,MJ in, 3.15\,MJ out; 2025-04-07 record: 8.6$\pm$0.45\,MJ out, target gain 4.13~\cite{llnl2025record}), and ITER official parameters. A subset (6 tasks) derives from \emph{verified local simulation assets} governed by the rules in Section~\ref{sec:verification}: a zero-dimensional Z-pinch dataset (10{,}944 rows; the identity $E=\tfrac12 mv^2$ holds in all 10{,}944 rows within $2\times10^{-4}$ relative deviation, two rows beyond the tighter $10^{-4}$ band only at rounding level), a screen-hydrogenic average-atom opacity model whose forced-single-core rerun reproduces all outputs byte-identically, and a FLASH~4.8 magnetohydrodynamic 15\,MA aluminum-liner Z-pinch setup used only for parameter/method questions because its physical outputs have not been independently rerun.

\textbf{Dual language track.} Tasks sourced from English-language material---15 problems built verbatim from real GitHub issues (DESC, OpenMC, WarpX, TORAX, BOUT++, FreeGS)---are \emph{English-canonical}: their canonical \texttt{question} is the original English issue text, with the Chinese version preserved in \texttt{question\_zh}/\texttt{options\_zh}. All other tasks are Chinese-canonical and additionally carry a first-class \texttt{question\_en} (and \texttt{options\_en} for choices), faithful translations checked by an automated physics-number fidelity gate (all hard numbers---decimals, scientific notation, $\geq$4-digit integers---must match character-exactly between the two languages; 0 mismatches at release). The harness evaluates either language via \texttt{--lang en|zh}; graders are language-agnostic (numbers, symbols, option letters, patches).

\begin{examplebox}{Box 1: Two verified tasks (abridged)}
\small
\textbf{Numeric (icf-ignition-gain-0003).} \emph{DT fusion releases 17.6\,MeV per reaction. NIF's first ignition shot yielded 3.15\,MJ; how many DT reactions occurred?} Answer: $1.117\times10^{18}$; \texttt{recompute\_expr} evaluates $3.15\times10^{6}/(17.6\times10^{6}\cdot1.602176634\times10^{-19})$ in a sandboxed math namespace and must agree within half the grading tolerance.\\[3pt]
\textbf{Derivation (icf-eos-opacity-0004).} \emph{Derive the Rosseland mean opacity in the two-group approximation: group opacities $k_1, k_2$ with weights $w_1, w_2$ (proportional to $\partial B_\nu/\partial T$).} Standard symbolic answer $(w_1+w_2)/(w_1/k_1+w_2/k_2)$ in the declared alphabet $\{w_1,w_2,k_1,k_2\}$; graded deterministically by SymPy symbolic equivalence (simplify-to-zero, proportionality for scaling forms, or sampled numeric equivalence).
\end{examplebox}

\begin{figure}[t]
\centering
\includegraphics[width=\linewidth]{fig_topic_dist.pdf}
\caption{Verified-split composition by topic and task type (v0.2, 688 tasks). ICF subdomains (blue) account for 50.9\%.}
\label{fig:topics}
\end{figure}
